\documentclass[a4paper,12pt]{article}

% Pacchetti utili
\usepackage[italian]{babel}
\usepackage[T1]{fontenc}
\usepackage[utf8]{inputenc}
\usepackage{lmodern}
\usepackage{geometry}
\usepackage{setspace}
\usepackage{graphicx}
\usepackage{titling}
\usepackage{hyperref}
\usepackage{booktabs}
\usepackage{amsmath, amssymb}
\usepackage[most]{tcolorbox}
\usepackage{tikz}

\geometry{margin=2.5cm}
\onehalfspacing

\begin{document}
	
	\begin{titlepage}
		\centering
		
		% Intestazione Università
		{\scshape\LARGE Università di Pisa \par}
		\vspace{0.5cm}
		
		% Titolo corso
		{\Huge\bfseries Calcolo Numerico - MATLAB\par}
		\vspace{0.5cm}
		{\Large Anno Accademico 2025--2026 \par}
		
		\vfill
		
		% Autore
		{\Large\textbf{Pietro Bacigalupi} \par}
		\vspace{0.5cm}
		
		\vfill
		
		% Data
		{\large \today\par}
		
	\end{titlepage}
	
	\tableofcontents
	\newpage
	\section*{Comandi MATLAB per l'Algebra Lineare}
	
	\subsection*{Creazione e Manipolazione di Matrici}
	\begin{itemize}
		\item \textbf{\texttt{A = [1 2 3; 4 5 6]}}: Il punto e virgola (\texttt{;}) serve per separare le righe durante la creazione di una matrice [cite: 442-443].
		\item \textbf{\texttt{A.'}} (Trasposta semplice): Questo operatore calcola la trasposta classica della matrice $A$, scambiando le righe con le colonne [cite: 440-441].
		\item \textbf{\texttt{A'}} (Trasposta coniugata / Hermitiana): Calcola la trasposta di $A$ e ne fa il complesso coniugato [cite: 440-441].
		\item \textbf{\texttt{eye(n)}}: Genera la matrice identità di ordine $n$[cite: 496].
		\item \textbf{\texttt{diag(v)}}: Crea una matrice con elementi uguali a 0 tranne che sulla diagonale principale, dove inserisce gli elementi del vettore $v$[cite: 450].
		\item \textbf{\texttt{diag(A)}}: Estrae i centri (diagonale principale) dalla matrice $A$[cite: 964].
		\item \textbf{\texttt{triu(A)}}: Estrae la parte triangolare superiore di una matrice data[cite: 452].
		\item \textbf{\texttt{tril(A)}}: Estrae la parte triangolare inferiore di una matrice data[cite: 452].
		\item \textbf{\texttt{A(2:n, [1:j-1, j+1:n])}}: Esclude la prima riga e la colonna $j$ usando gli indici, utile per lo sviluppo di Laplace[cite: 546].
		\item \textbf{\texttt{A(p, p)}}: Esegue la permutazione definendo il vettore $p$[cite: 1108].
	\end{itemize}
	
	\subsection*{Operazioni Algebriche}
	\begin{itemize}
		\item \textbf{\texttt{det(A)}}: Calcola il determinante della matrice quadrata $A$[cite: 473].
		\item \textbf{\texttt{inv(A)}}: Calcola l'inversa della matrice $A$[cite: 496].
		\item \textbf{\texttt{abs(B)}}: Prende la matrice con entrate i valori assoluti di $B$ [cite: 966-967].
		\item \textbf{\texttt{sum(C\_abs, 2)}}: Somma gli elementi sulle righe per ottenere i raggi[cite: 970].
	\end{itemize}
	
	\subsection*{Funzioni Grafiche e Vettoriali}
	\begin{itemize}
		\item \textbf{\texttt{linspace(-5, 10, 10000)}}: Genera punti equispaziati, ad esempio 10.000 punti tra -5 e 10[cite: 955].
		\item \textbf{\texttt{sin(t)}} e \textbf{\texttt{cos(t)}}: Applicano le funzioni seno e coseno per le parametrizzazioni [cite: 956, 960-961].
		\item \textbf{\texttt{plot(X, Y, 'r-o')}}: Traccia un grafico collegando i vettori delle ascisse e delle ordinate, dove \texttt{'r-o'} o \texttt{'g'} definiscono lo stile visivo [cite: 952-953, 957-962].
	\end{itemize}
	
\end{document}
